\documentclass[11pt,a4paper]{article}
\usepackage[margin=1in]{geometry}
\usepackage{hyperref}
\usepackage{xcolor}
\usepackage{listings}
\usepackage{booktabs}
\usepackage{enumitem}
\usepackage{fancyhdr}

\hypersetup{colorlinks=true,linkcolor=blue,urlcolor=blue}

\lstset{
  basicstyle=\ttfamily\small,
  backgroundcolor=\color{gray!10},
  frame=single,
  framerule=0.5pt,
  rulecolor=\color{gray!50},
  breaklines=true,
  breakatwhitespace=true,
  columns=fullflexible,
  keepspaces=true,
  showstringspaces=false,
  xleftmargin=4pt,
  xrightmargin=4pt,
  aboveskip=8pt,
  belowskip=8pt,
}

\pagestyle{fancy}
\fancyhf{}
\fancyhead[L]{\texttt{sheLLaMa}}
\fancyhead[R]{Cloud Fallback Setup Guide}
\fancyfoot[C]{\thepage}

\title{\textbf{Cloud Fallback Setup Guide}\\[4pt]\large OpenRouter \& LiteLLM for sheLLaMa}
\author{}
\date{}

\begin{document}
\maketitle
\thispagestyle{fancy}

\section{Overview}

sheLLaMa supports cloud LLM fallback when local Ollama models produce low-quality output. Two options are available:

\begin{itemize}
  \item \textbf{OpenRouter} --- cloud service at openrouter.ai, access to Claude/GPT-4/Llama/Gemini
  \item \textbf{LiteLLM} --- self-hosted proxy, runs on your network, routes to any LLM provider
\end{itemize}

When enabled, sheLLaMa automatically detects poor responses from the local model and re-sends the request through the configured fallback.

\section{Option 1: LiteLLM (Self-Hosted, Recommended)}

LiteLLM runs on your network as an OpenAI-compatible proxy. It can route to Ollama, vLLM, HuggingFace, or cloud providers.

\subsection{Install LiteLLM}

On a server in your network (can be the frontend or a dedicated host):

\begin{lstlisting}[language=bash]
pip install litellm[proxy]
\end{lstlisting}

\subsection{Configure LiteLLM}

Create a config file:

\begin{lstlisting}[language=bash]
cat > /etc/litellm/config.yaml << 'EOF'
model_list:
  - model_name: fallback
    litellm_params:
      model: ollama/qwen2.5-coder:14b
      api_base: http://192.168.1.218:11434
  - model_name: fallback
    litellm_params:
      model: ollama/qwen2.5-coder:14b
      api_base: http://192.168.1.219:11434
EOF
\end{lstlisting}

This routes the \texttt{fallback} model to your larger backends via Ollama directly. You can also add cloud models:

\begin{lstlisting}
  - model_name: cloud
    litellm_params:
      model: anthropic/claude-3.5-sonnet
      api_key: sk-ant-your-key-here
\end{lstlisting}

\subsection{Start LiteLLM}

\begin{lstlisting}[language=bash]
litellm --config /etc/litellm/config.yaml --port 4000
\end{lstlisting}

Or as a systemd service:

\begin{lstlisting}[language=bash]
cat > /etc/systemd/system/litellm.service << 'EOF'
[Unit]
Description=LiteLLM Proxy
After=network.target

[Service]
Type=simple
ExecStart=/usr/local/bin/litellm --config /etc/litellm/config.yaml --port 4000
Restart=always
RestartSec=3

[Install]
WantedBy=multi-user.target
EOF

systemctl daemon-reload
systemctl enable --now litellm
\end{lstlisting}

\subsection{Configure sheLLaMa to Use LiteLLM}

On each backend server:

\begin{lstlisting}[language=bash]
ssh root@<backend-ip>
sudo systemctl edit ansible-ollama
\end{lstlisting}

Add:

\begin{lstlisting}
[Service]
Environment="OPENROUTER_API_KEY=sk-anything"
Environment="OPENROUTER_MODEL=fallback"
Environment="OPENROUTER_URL=http://<litellm-host>:4000/v1/chat/completions"
Environment="USE_CLOUD_FALLBACK=true"
\end{lstlisting}

Note: LiteLLM requires an API key header but doesn't validate it by default --- any non-empty value works.

Restart:

\begin{lstlisting}[language=bash]
sudo systemctl restart ansible-ollama
\end{lstlisting}

\subsection{Configure All Backends at Once}

\begin{lstlisting}[language=bash]
for host in 192.168.1.230 192.168.1.233 192.168.1.218 192.168.1.219; do
  echo "--- $host ---"
  ssh root@$host "mkdir -p /etc/systemd/system/ansible-ollama.service.d && \
    cat > /etc/systemd/system/ansible-ollama.service.d/override.conf << 'EOF'
[Service]
Environment=\"OPENROUTER_API_KEY=sk-local\"
Environment=\"OPENROUTER_MODEL=fallback\"
Environment=\"OPENROUTER_URL=http://192.168.1.229:4000/v1/chat/completions\"
Environment=\"USE_CLOUD_FALLBACK=true\"
EOF
    systemctl daemon-reload && systemctl restart ansible-ollama"
done
\end{lstlisting}

\subsection{Verify LiteLLM}

\begin{lstlisting}[language=bash]
curl http://<litellm-host>:4000/v1/models
curl -X POST http://<litellm-host>:4000/v1/chat/completions \
  -H "Authorization: Bearer sk-local" \
  -H "Content-Type: application/json" \
  -d '{"model": "fallback", "messages": [{"role": "user", "content": "hello"}]}'
\end{lstlisting}

\section{Option 2: OpenRouter (Cloud)}

\begin{enumerate}
  \item Go to \url{https://openrouter.ai}
  \item Click \textbf{Sign Up} and create an account (Google/GitHub OAuth or email)
  \item Navigate to \textbf{Keys} at \url{https://openrouter.ai/keys}
  \item Click \textbf{Create Key}
  \item Name it (e.g., \texttt{shellama}) and copy the key --- it starts with \texttt{sk-or-v1-...}
\end{enumerate}

\section{Add Credits}

\begin{enumerate}
  \item Go to \textbf{Credits} at \url{https://openrouter.ai/credits}
  \item Add funds --- \$5 is enough for significant usage
  \item Optionally set a usage limit to avoid surprises
\end{enumerate}

\section{Choose a Model}

Browse available models at \url{https://openrouter.ai/models}.

\begin{table}[h]
\centering
\begin{tabular}{llll}
\toprule
\textbf{Model} & \textbf{ID} & \textbf{Cost (per 1M tokens)} & \textbf{Notes} \\
\midrule
Claude 3.5 Sonnet & \texttt{anthropic/claude-3.5-sonnet} & \textasciitilde\$3/\$15 & Default, high quality \\
GPT-4o & \texttt{openai/gpt-4o} & \textasciitilde\$2.50/\$10 & Fast, good quality \\
Llama 3 70B & \texttt{meta-llama/llama-3-70b-instruct} & \textasciitilde\$0.60/\$0.80 & Open source, cheap \\
Gemini Pro 1.5 & \texttt{google/gemini-pro-1.5} & \textasciitilde\$1.25/\$5 & Large context window \\
\bottomrule
\end{tabular}
\end{table}

\section{Configure sheLLaMa Backends}

The cloud fallback is configured on each \textbf{backend} server (not the frontend). You need to set three environment variables in the systemd service.

\subsection{Option A: Using systemd override (recommended)}

On each backend server:

\begin{lstlisting}[language=bash]
ssh root@<backend-ip>
sudo systemctl edit ansible-ollama
\end{lstlisting}

This opens an editor. Add the following:

\begin{lstlisting}
[Service]
Environment="OPENROUTER_API_KEY=sk-or-v1-your-key-here"
Environment="OPENROUTER_MODEL=anthropic/claude-3.5-sonnet"
Environment="USE_CLOUD_FALLBACK=true"
\end{lstlisting}

Save and restart:

\begin{lstlisting}[language=bash]
sudo systemctl restart ansible-ollama
\end{lstlisting}

\subsection{Option B: Edit the service file directly}

\begin{lstlisting}[language=bash]
ssh root@<backend-ip>
sudo nano /etc/systemd/system/ansible-ollama.service
\end{lstlisting}

Update the \texttt{Environment} lines:

\begin{lstlisting}
[Service]
Environment="OPENROUTER_API_KEY=sk-or-v1-your-key-here"
Environment="OPENROUTER_MODEL=anthropic/claude-3.5-sonnet"
Environment="USE_CLOUD_FALLBACK=true"
\end{lstlisting}

Reload and restart:

\begin{lstlisting}[language=bash]
sudo systemctl daemon-reload
sudo systemctl restart ansible-ollama
\end{lstlisting}

\subsection{Configuring All Backends at Once}

\begin{lstlisting}[language=bash]
for host in 192.168.1.230 192.168.1.233 192.168.1.218 192.168.1.219; do
  echo "--- $host ---"
  ssh root@$host "mkdir -p /etc/systemd/system/ansible-ollama.service.d && \
    cat > /etc/systemd/system/ansible-ollama.service.d/override.conf << 'EOF'
[Service]
Environment=\"OPENROUTER_API_KEY=sk-or-v1-your-key-here\"
Environment=\"OPENROUTER_MODEL=anthropic/claude-3.5-sonnet\"
Environment=\"USE_CLOUD_FALLBACK=true\"
EOF
    systemctl daemon-reload && systemctl restart ansible-ollama"
done
\end{lstlisting}

\section{Verify}

Check that the service is running with the new config:

\begin{lstlisting}[language=bash]
ssh root@<backend-ip>
sudo systemctl status ansible-ollama
\end{lstlisting}

Test the fallback by sending a request that would produce a poor local response:

\begin{lstlisting}[language=bash]
curl -X POST http://<backend-ip>:5000/generate-code \
  -H "Content-Type: application/json" \
  -d '{"description": "complex distributed system", "model": "qwen2.5-coder:7b"}'
\end{lstlisting}

If the local model's response is empty or an error, the response JSON will include:

\begin{lstlisting}
{
  "cloud_fallback": true,
  "cloud_model": "anthropic/claude-3.5-sonnet"
}
\end{lstlisting}

\section{Disable Cloud Fallback}

To go back to fully offline operation:

\begin{lstlisting}[language=bash]
ssh root@<backend-ip>
sudo systemctl edit ansible-ollama
\end{lstlisting}

Change to:

\begin{lstlisting}
[Service]
Environment="USE_CLOUD_FALLBACK=false"
\end{lstlisting}

Restart:

\begin{lstlisting}[language=bash]
sudo systemctl restart ansible-ollama
\end{lstlisting}

\section{How Fallback Works}

\begin{enumerate}
  \item Every request is first processed by the local Ollama model
  \item The response is checked for quality:
    \begin{itemize}[nosep]
      \item Empty playbook or code output $\rightarrow$ triggers fallback
      \item Error from Ollama $\rightarrow$ triggers fallback
      \item Successful non-empty response $\rightarrow$ returned as-is (no cloud call)
    \end{itemize}
  \item If fallback triggers, the same prompt is sent to OpenRouter
  \item The cloud response replaces the local response
  \item The response is tagged with \texttt{cloud\_fallback: true} so clients know
\end{enumerate}

\section{Environment Variables Reference}

\begin{table}[h]
\centering
\begin{tabular}{lll}
\toprule
\textbf{Variable} & \textbf{Default} & \textbf{Description} \\
\midrule
\texttt{OPENROUTER\_API\_KEY} & \textit{(empty)} & API key (OpenRouter or any string for LiteLLM) \\
\texttt{OPENROUTER\_MODEL} & \texttt{anthropic/claude-3.5-sonnet} & Model to use for fallback \\
\texttt{OPENROUTER\_URL} & \texttt{https://openrouter.ai/...} & API endpoint (change for LiteLLM) \\
\texttt{USE\_CLOUD\_FALLBACK} & \texttt{false} & Set to \texttt{true} to enable \\
\bottomrule
\end{tabular}
\end{table}

\section{Cost Management}

\begin{itemize}
  \item Monitor usage at \url{https://openrouter.ai/activity}
  \item Set spending limits at \url{https://openrouter.ai/credits}
  \item Use cheaper models (e.g., \texttt{meta-llama/llama-3-70b-instruct}) to reduce costs
  \item Fallback only triggers on poor local responses, so costs are minimal with a good local model
\end{itemize}

\end{document}
